% The one Technical Background file that describes the prediction model.
% Trimmed 2026-11-07 from thesis/technical_background/01_risk_models.tex, 02_segment_glmm.tex and
% the reserve 02_segment_glmm_long.tex (full drafts in git, commit e4eb857).
% 2026-11-13 (user: "use the pooled cloglog GLMM evaluation to write the technical section; style of
%   fpca.tex first half; flow between the lines"). WRITTEN: the GLMM part only (scope chosen with the
%   user; risk models, recalibration and the general bootstrap are still NOTES below, because the base
%   choice, recalibrated vs refitted SCORE2, is open). 4 displayed equations (eq:tech-glmm,
%   eq:tech-lr, eq:tech-agg-conditional, eq:tech-agg-marginal). Source: reports/Pooled cloglog GLMM
%   evaluation.md, on top of the retired outdated/thesis_retired/technical_background/02_segment_glmm.tex.
%   Argument changed from the retired draft: the random intercept is justified by valid standard errors
%   and a joint distribution, NOT by better prediction; "slope 0.79 vs 0.98" dropped (logit, ICC 0.30,
%   not the size here); cloglog is the grouped-time form of a PH frailty model; log exposure with a
%   free coefficient and the offset as the special case; variance weakly identified in small binary
%   clusters; LR test needs the random intercept to capture the dependence; the independence product
%   overstates the cluster risk. Generic wording on purpose ("unit", "cluster", "exposure"): the
%   segments, persons, log length and the nine segments are Methodology. Notation: H_i (summed
%   hazard, not S_i, taken by the shrinkage in 01b_sample_size.tex), q (block size, not k), f(b;
%   sigma^2) for the normal density (phi is the FPCA eigenfunction).
%   Not cited on purpose (individual studies, Lit. Review or Discussion): Lingsma 2013, Bouwmeester
%   2013, Falconieri 2020, TRIPOD-Cluster, the coronary cohort papers, Riley.
%   2026-11-13 (counterfactual rule): GEE cut from the opener. Removed: "Marginal models such as GEE
%     estimate the average effect across clusters and correct the SEs without describing the
%     correlation (liang1986gee)" and "for a new cluster, predictions about as accurate as a marginal
%     model once b is integrated out (pavlou2015marginal)". The real case for GLMM over GEE is the
%     joint distribution needed for the cluster risk (sec:tech-aggregation); GEE SEs survive a wrong
%     working correlation, GLMM SEs and the LR test do not. For the examiner and, if the artery-level
%     dependence is raised, one line in the Discussion. liang1986gee is now uncited.
%   CHECK: pavlou2015marginal is cited for "predictions about as accurate as a marginal model once
%     integrated" (report: opened via a summarising tool, logit simulations); confirm in the PDF.
%   CHECK: moineddin2007multilevel and austin2010clusters, same (poor variance estimates, small
%     binary clusters; regimes are 5-100 clusters, not ~1000 clusters of nine).
%   CHECK: prentice1978grouped, duchateau2008frailty (chapter and page), liu1994gausshermite,
%     pinheiro1995approximations, joe2008laplace were not read in full; the unread ones are cited for
%     standard results (grouped PH gives cloglog; shared frailty; adaptive quadrature; Laplace bias).
%   CHECK: "the test then rejects too often" and "the independence product overstates" are reasoning,
%     uncited; the report grades them as inference.
%   Pointed back to (labels): sec:tech-segment-glmm, sec:tech-glmm-likelihood, sec:tech-aggregation,
%     eq:tech-lr, sec:tech-nested (methodology/03_prediction_model.tex notes).
% 2026-11-13: added sec:tech-nested (nested random intercepts, eq:tech-nested, eq:tech-nested-icc) for
%   the three-level sensitivity analysis (person + artery within person; mapping is Methodology).
%   Notation: subgroup g(j) keeps y_ij/eta_ij; omega^2 (tau is the taper, k the coefficient index).
%   Boundary LR test cites self1987boundary, stram1994variance [TO VERIFY]. Latent variance pi^2/6
%   has a \todonote until a source is read. sec:tech-aggregation still assumes one intercept: fine
%   while the sensitivity analysis reports only variances, LR test and hazard ratios; extend it only
%   if person risk is computed under the nested model.
%   CHECK: file spelling is American (summarized, neighboring) against the British rule in CLAUDE.md;
%   new text follows the file ("unmodelled" is the one British form); sweep later.
% Added value (C, calibration, Brier, log loss, NRI/IDI, internal validation; was 03_added_value.tex)
%   is deliberately not written yet.
% Pointed back to by: content/methodology/03_prediction_model.tex (sec:method-models).
% Rule: general methods only, present tense; settings, covariates, software and node counts are
% Methodology.
%
% The section should contain, whatever the model:
% - Clinical risk prediction models: absolute risk of an outcome over a time horizon from
%   predictors, linear predictor mapped to a risk through a link (steyerberg2019cpm). Weights set
%   the ranking, the intercept the level.
% - Developed in the data versus a published score; SCORE2 (score2_2021) as the example. A published
%   score costs no parameters but its weights were chosen for its own outcome and population.
% - Recalibration (janssen2008updating, cox1958binary): in-the-large, logistic recalibration
%   (intercept + one slope on the published linear predictor), revision, refitting; each step costs
%   more parameters. Recovering the linear predictor from a published risk on the link scale.
% - Hazard closer: a recalibrated score keeps its relative weights, so a new block correlated with a
%   mis-weighted factor (e.g. age) can look informative only by correcting that weight. Does NOT
%   decide the open base choice (recalibrated SCORE2 vs refitted factors).
%
% - The bootstrap in general (added 2026-11-09): resampling the observations (here persons) with
%   replacement to see how an estimate would vary across samples; percentile intervals. Used twice
%   in the thesis: the person-level bootstrap for the intervals of the model's estimated measures
%   (Methods 03 notes; plan in ~/Desktop/Speciale/CGPS_statistics-/CLAUDE.md) and the stability
%   check of the FPCA components (sec:tech-fpca-components, Methods 07). Explain it once here so
%   the FPCA section can point back. CITATION WANTED: a classic source (e.g. Efron and Tibshirani
%   1993, An Introduction to the Bootstrap); not yet in refs.bib, add and check before citing.
%
% If the cloglog GLMM is confirmed, it should also contain:
% (Updated 2026-11-11 from reports/Pooled cloglog GLMM evaluation.md. Its verdict: keep the model,
% change the argument for it. No non-ML alternative was found that is clearly better.)
% - Clustered binary outcomes (segments within persons); correlated outcomes make naive standard
%   errors too small; GLMM adds a random person effect (statacorp2025mecloglog). Practice: only
%   5 of 15 intracoronary imaging studies with several observations per patient handled clustering
%   (Lingsma 2013 EuroIntervention, NOT IN refs.bib); coronary cohort papers mostly use GEE or
%   marginal Cox with robust variance, so say why a GLMM rather than GEE (see "why a GLMM" below).
% - WHY A GLMM (the justification to write; replaces any "predicts better" claim): a random
%   intercept does NOT predict better for a new person. Integrated predictions tie with GEE and
%   standard logistic (calibration slopes 0.97-0.98, pavlou2015marginal); b = 0 is worse
%   (Bouwmeester 2013, NOT IN refs.bib). The model earns its place for two other reasons: (1)
%   honest standard errors, so a valid LR test of the geometry block; (2) a joint model of the
%   segments, which is what lets segment risks be combined into a person risk. GEE gives only
%   marginal segment risks, and the independence product 1 - prod(1 - p) overstates the person
%   risk under positive within-person correlation.
% - Cloglog link: proportional hazards over the window give hazard ratios independent of window
%   length (prentice1978grouped); logit gives odds ratios, close at low risk. Say honestly that fit
%   cannot tell the links apart at segment level; the reason for cloglog is interpretive (hazard
%   ratios, segment hazards add to a person hazard). The link does matter for the person risk,
%   which is not small (check the event rate, see below).
% - Survival reading (use as justification): with two scans each segment has one interval, so an
%   interval-censored shared-frailty PH model collapses to this GLMM: the two-scan, discrete-time
%   counterpart of the shared-frailty Cox model that lesion-level studies already accept. Log
%   follow-up as an offset = constant baseline hazard; free coefficient = Weibull baseline
%   (standard result, VERIFY in prentice1978grouped before writing). PH cannot be tested with one
%   follow-up scan; say so.
% - Exposure such as segment length: log length as an offset (coefficient 1) assumes hazard
%   proportional to length, with no coronary precedent found. Present the free coefficient as the
%   default (1 df, interval reported against 1) and the offset as the special case. Curvature and
%   taper are computed per segment and depend on length, so a wrong offset can leak into the
%   tested block.
% - Random intercept: shared frailty multiplying all segment hazards (duchateau2008frailty, page
%   still to verify), summarised by sigma^2; conditional versus marginal effects
%   (pavlou2015marginal); median hazard ratio (austin2017mhr). Weak point to state: sigma^2 is
%   informed almost only by persons with two or more positive segments, and small binary clusters
%   bias variance estimates (Moineddin 2007, Austin 2010, NOT IN refs.bib; their regimes are 5-100
%   clusters, not ~1000 clusters of nine, so the size of the bias here is unknown). Robustness to
%   a non-normal random effect holds for slopes, not for variance components or the person risk.
% - Fitting and testing: integration over the random intercept by adaptive Gauss-Hermite quadrature
%   (liu1994gausshermite); Laplace bias in small binary clusters (joe2008laplace), so Laplace is
%   never the primary fit; LR test of a block of k coefficients against chi^2_k, both models on
%   the same rows and settings. The chi^2 reference is asymptotic; a parametric bootstrap of the
%   LR statistic checks it, but only under the fitted model (it does not protect against
%   dependence beyond one shared intercept; a person-cluster bootstrap or robust Wald test does).
% - Segment risk to person risk: given b_i the segment hazards add (S_i), and the person risk
%   integrates over b; b_i = 0 gives the median person and is miscalibrated. Do NOT quote
%   pavlou2015marginal's "slope 0.79 vs 0.98" as the expected size here: it is a logit result for
%   individual predictions at ICC 0.30. Use it for the direction only; the size is computed in the
%   cohort (Methods). The person risk depends on sigma^2 and on the assumed normal shape, which
%   segment-level fit and the C-statistic cannot see; only person-level calibration can. An effect
%   on one segment is diluted in S_i, so a rejected LR test need not change the person risk.
% - Sample size: Riley criteria (riley2019) assume independent observations and are silent on
%   clustering. Count persons and person-level events, never segment rows (three of the five
%   geometry features are person-level constants, so segment rows are pseudo-replicates for
%   them). Shrinkage of the pooled model is estimated by resampling, not asserted from the formula.
% - TRIPOD-Cluster (Debray 2023, NOT IN refs.bib) explicitly puts clustering within individuals
%   out of scope; cite it by analogy only: state how clustering was handled, which intercept the
%   predictions use, and keep every model-building step inside the resampling.
%
% Belongs in Methodology 03, not here (recorded so it is not lost; general method only above):
% - Base model: free coefficient on log length; integrated person risk stated explicitly, b = 0 and
%   the independence product not used; the b = 0 error computed in the cohort.
% - Decide which rows carry the three left-main features (now pLAD and pLCx, adjacent rows) before
%   the test is locked; put a person-level test of those three features beside the block test.
% - Before locking: tabulate the segment-level event rate and the number of persons with two or
%   more positive segments (they decide how well sigma^2 is identified). The research notes
%   assumed rare events; the earlier report gives ~452 of 1,076 persons with new plaque (~42%),
%   NOT verified against current data.
% - Prespecified sensitivity analyses, in order of value: (1) person-level cloglog on per-person
%   geometry summaries (benchmark; tests "who"); (2) working-independence cloglog GEE with a
%   robust 5-df Wald test plus person-cluster bootstrap; (3) parametric bootstrap of the LR
%   statistic; (4) gamma-frailty refit (closed-form person risk); (5) fold-wise sigma-hat from
%   several starts, boundary estimates reported; (6) logit refit and fixed-offset length; (7)
%   person means of curvature and taper (+2 df, between-within); (8) number of evaluable segments
%   as covariate, calibration by segment count; (9) observed vs model-implied plaque
%   co-occurrence in adjacent and non-adjacent segment pairs; (10) vessel-within-person intercept.
% - Dependence beyond one shared intercept (added 2026-11-11, from talk prep). The single b_i
%   gives every pair of segments the same latent dependence, but real dependence likely follows
%   the tree: parent-child along a vessel (proximal -> mid -> distal) and the two left-main
%   daughters (pLAD-pLCx), which also carry the same left-main feature values. If so, those rows
%   count as more independent evidence than they are and the LR test is anticonservative. Do NOT
%   add the vessel intercept (10) to the primary model: it is a second variance, informed only by
%   persons with >= 2 positives in the same vessel, so it will likely sit at the boundary; keep
%   it as a sensitivity fit. Check the assumption without spending parameters instead:
%   (a) Co-occurrence check (9): compare how often a pair of segments is both positive in the
%       data with how often the fitted model predicts it, separately for parent-child pairs, the
%       pLAD-pLCx siblings and pairs in different vessels. Uses b: the model-implied joint
%       probability is P(Y_j = 1, Y_k = 1) = integral of p_j(b) p_k(b) N(b; 0, sigma^2) db, the
%       segments being independent only given b. If observed and predicted agree, the single
%       sigma^2 is adequate.
%   (b) Person-cluster bootstrap or GEE robust Wald test (2): valid for the geometry block
%       whatever the true dependence (tree-shaped, sibling or other), without modelling it. If
%       they agree with the LR test, the result holds. The cluster bootstrap resamples whole
%       persons and refits the GLMM, so it still uses b. The GEE has no b: it is a marginal
%       (population-averaged) model with working independence and sandwich standard errors, so
%       its coefficients are on a different (attenuated) scale from the GLMM's conditional ones;
%       compare the tests (geometry block = 0), not the coefficient sizes.
% - Optional: a short simulation at the fitted base-model parameters (sigma-hat = 0 rate, type I
%   error of LR and robust Wald for a 5-df block, person-level calibration under normal and
%   non-normal frailty). No published simulation covers this data shape.
% - Rejected alternatives, one line each in the Discussion: penalised (ridge/lasso/Firth) GLMM
%   (unreliable at small effective n, loses the LR test; Riley 2021), Bayesian GLMM (equal at
%   best), beta-binomial (cannot hold segment covariates), conditional logistic (drops persons
%   without events and all person-level features), random slopes and spatial effects (not
%   estimable with nine rows).
% - Examiner checklist: GLMM vs GEE; cloglog vs logit; length as exposure and whether fixed;
%   conditional vs marginal and which prediction is evaluated; the aggregation assumption;
%   resampling by person including every estimated quantity; quadrature settings; how a sample-size
%   criterion for independent persons was applied to a model fitted on segment rows.
% - CHECK BEFORE CITING: every number above was relayed by a summarising tool; check against the
%   PDFs. Full list of abstract-only and snippet-only sources in the report's last section.
% 2026-11-13 (user: "only the most important, 2-3 pages; human phrasing after the Aida style guide,
%   not a math book; easy on the first read; concise"). Rewritten from about 1,500 to about 900 prose
%   words, 6 to 5 displayed equations, 4 to 3 subsections. Copy before this pass: scratchpad
%   01_prediction_model_before_short.tex (not in the repo). Claims, cite keys and kept labels unchanged.
%   - FOLDED (user): sec:tech-nested is now one paragraph inside sec:tech-segment-glmm (label kept on
%     the paragraph, so a reference gives that subsection's number). eq:tech-nested kept.
%   - CUT: eq:tech-nested-icc (latent correlations rho_same, rho_different) and its pi^2/6 todonote
%     (no source read); nothing referenced the label. Restore from git if the Methodology needs ICCs.
%   - CUT: "the logit odds ratios change with the baseline risk"; the cost of nested quadrature
%     ("grows with the square of the number of nodes"); the dilution algebra (H_i multiplied by
%     (1 - w) + wh, w the unit's share), now said in words; "attenuated toward zero" -> "smaller".
%   - Boundary test reduced to one sentence (self1987boundary, stram1994variance [TO VERIFY]).
%   - Spelling now British in the prose (generalised, summarised, maximised); heading kept as
%     "Generalized" until the whole-thesis spelling sweep.
% 2026-11-13 (second pass, user: "look at how Aida explains mathematical components; not static and
%   academic; human-written; go through the entire text"). Aida's pattern applied: assumptions in
%   words first ("Under these assumptions, ..."), then the equation, then the symbols in a "where"
%   clause after it; no symbol defined before the reader knows why it is needed (the opening "Let
%   y_ij be 1 if ... and let eta_ij be ..." is gone). Each step of the chain says what it adds ("So
%   far the units are treated as independent"). NEW: one unnumbered display for the model without
%   the cluster effect, so eq:tech-glmm reads as one added term. y_ij no longer named in the text;
%   A_i is defined in words in sec:tech-aggregation. Claims, cites and labels unchanged.
% 2026-11-13 (third pass, user: "make the technical background more human"): light touch only, the
%   section was already rewritten for tone twice today. "Its main appeal is the reading of" -> "that
%   the coefficients are easy to read:"; "fit nearly alike" -> "give nearly the same fit"; "Since"
%   opener removed; ambiguous "It ranks" -> "This shortcut ranks"; dilution closer reworded.
% 2026-11-13 (user): first display rewritten so the hazard ratio is visible in it: Lambda_ij (hazard
%   accumulated over the window) = exp(eta_ij), eta_ij = beta_0 + sum_k beta_k x_ijk (was x^T beta,
%   with beta_k never defined). Link written out; why window-independence matters added. Later
%   displays unchanged: exp(eta_ij + b_i) is Lambda with the cluster effect, and H_i in
%   sec:tech-aggregation is the sum of the Lambda_ij.
% 2026-11-13 (user: better openings, collapse sections): subsection "The Complementary Log-Log Model"
%   removed, so the section runs from the clustering problem straight into the model; its label
%   sec:tech-segment-glmm now sits on the paragraph (nothing outside this file references it).
%   Openers: "Some units have more room ..." -> "Units also differ in how much opportunity ...";
%   "So far the units are treated as independent." -> "The model above still treats every unit on
%   its own. To let the units of a cluster share their risk, ..." (", shared by all units of
%   cluster i," inserted clause removed). Two subsections remain: Fitting and Testing; From Unit
%   Risk to Cluster Risk.
% 2026-11-13 (user, /plan: "The Complementary Log-Log Model more concise, not as deep where not
%   necessary"): subsection is back in the text and shortened by about a third. CUT: "This matters when
%   the windows differ in length"; the full median-hazard-ratio definition; "b_i not estimated one by
%   one" (merged with sigma^2 = 0); "units in the same subgroup share both effects" (follows from
%   eq:tech-nested); "hazard grows as E^gamma". Claims, cites, equations and labels unchanged.
% 2026-11-13 (user: "infer aida styleguide over the section"; outdated/thesis_retired/
%   style_guide_technical_aida.md + style_guide_aida.md): assumptions named before the first display
%   ("Under the assumptions that ..."); symbols in "where" clauses (sigma^2, omega^2 back there);
%   "event rate, or hazard" synonym pattern; "A key characteristic of this link"; random effect and
%   nested effect open on the limitation they address ("While ...", "A single random effect makes
%   ..."); inserted glosses removed (", the hazard, which", ", an offset,"); participial tail on
%   exposure; sigma^2 hazard closer split into fact + "therefore". Slightly longer than the concise
%   pass, still shorter than the original. Claims, cites, equations and labels unchanged.
% 2026-11-13 (user: "the prediction model cannot be longer than Aida's longest section"): Aida's
%   longest is 3.2 U-Net + 3.2.1 nnU-Net, about 735 words incl. equation fragments (3.3 + 3.3.1
%   about 730). Whole section cut from about 1,260 to under 735 on the same basis. Labels, cites
%   and numbered equations kept. [2026-11-13 later, user: "more academic, keep tone": cloglog opening
%   reworded; assumption named (proportional hazards); "hazard accumulated" -> "cumulative hazard";
%   link written as log{-log(1-p)} = eta; "A key property of this link". Claims and cites unchanged.
%   Then (user: "more human and direct, optimise flow"): plain opener ("either has an event ... or it
%   does not"); assumption as a colon statement; "Rearranging gives"; "Its main advantage";
%   exposure introduced by its reason ("Units also differ in how much opportunity ...").
%   Then (user: "concise, Aida style, no repetition"): "A key characteristic of this link" (Aida 3.3);
%   HR defined per unit increase; exposure via "commonly referred to as"; gamma's role stated as
%   hazard ~ E^gamma; example "such as the size of the unit" dropped (inserted gloss, and the chapter
%   stays generic). No re-mention of cloglog name, PH or window definition.]
%   CUT: logit comparison (counterfactual rule; one line for the
%   Discussion: logit gives odds ratios, near-identical fit when events are rare); exposure as
%   E^gamma; sigma^2 = 0 'no clustering'; 'effects averaged over clusters are smaller'; nested
%   'informed only by subgroups with > 1 event, can come out as zero'; 'more nodes ... closer';
%   'same settings, so the difference comes from the block'; 'nodes raised until estimates stop
%   changing' (a setting, Methodology); nested quadrature sentence; 'integral
%   stays one-dimensional'; 'only calibration can tell them apart' (folded); dilution closer
%   shortened to one sentence; eq:tech-nested display made inline (label dropped, nothing cited it;
%   sec:tech-nested kept). Backup of the pre-cut file: scratchpad 01_prediction_model_before_735.tex.
\section{Generalized Linear Mixed Model}
\label{sec:tech-model}

A generalised linear mixed model (GLMM) is a regression model designed for outcomes observed in clusters, where several units belong to the same cluster.
Units of one cluster share unmeasured traits, so their outcomes tend to occur together, and ignoring this gives standard errors that are too small.
A GLMM adds a random effect for each cluster, keeping the standard errors valid and modelling the dependence within a cluster \cite{statacorp2025mecloglog}.

\subsection{The Complementary Log-Log Model}
\label{sec:tech-segment-glmm}
Each unit is followed over an observation window, and the outcome is whether it has an event.
To link this outcome to the covariates, the model assumes proportional hazards: each covariate scales the event rate by a constant factor throughout the window.
The probability that unit $j$ of cluster $i$ has an event is then given by 
\[
  p_{ij} = 1 - \exp\!\left(-\Lambda_{ij}\right),
  \qquad
  \Lambda_{ij} = \exp(\eta_{ij}),
  \qquad
  \eta_{ij} = \beta_0 + \sum_{k} \beta_k\, x_{ijk},
\]
where $\Lambda_{ij}$ is the hazard summed over the whole window, known as the cumulative hazard, and $\beta_k$ is the coefficient of covariate $x_{ijk}$.
Rearranging shows that the model is linear in the covariates on a transformed scale,
\[
  \log\{-\log(1 - p_{ij})\} = \eta_{ij},
\]
where the left-hand side is the complementary log-log (cloglog) of the event probability.
Increasing $x_{ijk}$ by one unit multiplies the hazard by $\exp(\beta_k)$, which is therefore the hazard ratio of covariate $k$.
The factor applies equally at every moment of the window, so the hazard ratio does not depend on how long the window is \cite{prentice1978grouped}.
Units may also differ in how much opportunity they have for an event, for example because some units are larger than others.
This opportunity is called the exposure, $E_{ij}$, and it enters $\eta_{ij}$ as an additional term $\gamma \log E_{ij}$.
The hazard then grows in proportion to $E_{ij}^{\gamma}$, so that with $\gamma = 1$ doubling the exposure doubles the hazard.
When $\gamma$ is fixed at 1 instead of being estimated, the term is called an offset.

While this model treats every unit on its own, the units of a cluster share part of their risk.
The GLMM captures this by adding a random effect $b_i$ to every unit in cluster $i$,
\begin{equation}
  p_{ij}(b_i) = 1 - \exp\!\left\{-\exp(\eta_{ij} + b_i)\right\},
  \qquad
  b_i \sim \mathcal{N}(0, \sigma^2),
  \label{eq:tech-glmm}
\end{equation}
where $\sigma^2$ is the variance of the cluster effects, often reported as a median hazard ratio \cite{austin2017mhr}.
Given $b_i$ the units are assumed independent, and $e^{b_i}$ multiplies the hazard of every unit in the cluster by the same factor, known as a shared frailty \cite{duchateau2008frailty}.
The coefficients are therefore conditional, comparing two units of the same cluster \cite{pavlou2015marginal}.
Most of the information about $\sigma^2$ comes from clusters with more than one event, so it is estimated poorly when clusters are small and events are few \cite{moineddin2007multilevel, austin2010clusters}.

\label{sec:tech-nested}%
When the units fall into subgroups that share more than the cluster as a whole, a nested random intercept adds a second effect $c_{ig(j)} \sim \mathcal{N}(0, \omega^2)$ for the subgroup $g(j)$ of unit $j$ next to $b_i$ in Eq.~\eqref{eq:tech-glmm}.
Setting $\omega^2 = 0$ gives back the single intercept, and $\omega^2$ is even harder to estimate than $\sigma^2$.

\subsection{Fitting and Testing}
\label{sec:tech-glmm-likelihood}

The cluster effects are not observed, so the likelihood averages over $b_i$, an integral with no closed form.
Adaptive Gauss--Hermite quadrature approximates it by a weighted sum over a few nodes placed around the most likely value of $b_i$ \cite{liu1994gausshermite, pinheiro1995approximations}.
With a single node it reduces to the Laplace approximation, which is biased for binary outcomes in small clusters \cite{joe2008laplace, statacorp2025mecloglog}.

To test whether a block of $q$ covariates adds anything, the model is fitted with and without the block on the same rows,
\begin{equation}
  \mathrm{LR} = 2\left(\ell_1 - \ell_0\right) \;\overset{H_0}{\sim}\; \chi^2_q ,
  \label{eq:tech-lr}
\end{equation}
where $\ell_1$ and $\ell_0$ are the maximised log-likelihoods.
The $\chi^2$ reference holds only if the random effects capture the dependence within a cluster, and otherwise the test rejects too often.
For $\omega^2 = 0$ the p-value of the $\chi^2_1$ test is halved, as a variance cannot be negative \cite{self1987boundary, stram1994variance}.

\subsection{From Unit Risk to Cluster Risk}
\label{sec:tech-aggregation}

The question of interest can be whether any unit of a cluster has an event, written $A_i = 1$.
Given $b_i$ the units are independent, and under the cloglog model the product of their probabilities collapses to a single exponential,
\begin{equation}
  P(A_i = 0 \mid b_i) = \prod_{j=1}^{n_i} \exp\!\left(-e^{\eta_{ij} + b_i}\right)
  = \exp\!\left(-e^{b_i} H_i\right),
  \qquad
  H_i = \sum_{j=1}^{n_i} e^{\eta_{ij}} ,
  \label{eq:tech-agg-conditional}
\end{equation}
where $n_i$ is the number of units and $H_i$ is their summed hazard.
For a new cluster $b_i$ is unknown, so the risk is averaged over its distribution,
\begin{equation}
  P_i = P(A_i = 1) = 1 - \int \exp\!\left(-e^{b} H_i\right) f(b;\sigma^2)\,\mathrm{d}b ,
  \label{eq:tech-agg-marginal}
\end{equation}
where $f(b;\sigma^2)$ is the normal density with mean 0 and variance $\sigma^2$.

Setting $b_i = 0$ instead gives the risk of the median cluster, which ranks the clusters identically but is miscalibrated \cite{pavlou2015marginal, skrondal2009prediction}.
Multiplying the unit risks as if the units were independent overstates the cluster risk when outcomes occur together.
It also depends on $\sigma^2$ and the normal shape of $b_i$, so it is judged by calibration at the cluster level.
An effect on a single unit is diluted in $H_i$, so a block that is significant in Eq.~\eqref{eq:tech-lr} may barely change the cluster risk.
